\subsection{\label{ref:eq_optimizer}EQ Optimizer}
EQ Optimizer reduces the number of bands used by the current
\setting{Equalizer} settings (see \reference{ref:EQ}) while keeping the
frequency response as close as possible to the original.  Each enabled EQ
band costs CPU time while music plays, so complex EQ settings, such as those
generated by headphone correction tools, can be replaced by settings with
fewer bands that sound almost the same and use less battery.

When started, the plugin reads the current EQ settings and fits them with
every smaller number of bands, down to one.  At least two EQ bands must be
enabled.  The fitting can take several seconds on slower players.  It then
shows a plot of the original frequency response (solid line) and the fitted
response (\opt{lcd_color}{orange line}\nopt{lcd_color}{dashed line}) from
20~Hz to 20~kHz.  A strip along the bottom of the screen shows the difference
between the two curves, magnified; the number at its left edge is the
difference in dB that fills the strip.  The top of the screen shows the
number of bands used by the fit, and its RMS and maximum error in dB compared
to the original.  On players with a large enough screen, the estimated CPU
load of the EQ and the saving compared to the original settings is also shown.

The plugin starts on the fit with the fewest bands that has an RMS error of
no more than 0.25~dB and a maximum error of no more than 1~dB.

\begin{btnmap}
    \opt{scrollwheel}{\PluginScrollBack{} / \PluginScrollFwd{} or}
    \PluginLeft{} / \PluginRight
        \opt{HAVEREMOTEKEYMAP}{& \PluginRCLeft{} / \PluginRCRight}
        & Use fewer / more bands\\

    \PluginSelect
        \opt{HAVEREMOTEKEYMAP}{& \PluginRCSelect}
        & Open the menu\\

    \nopt{IPOD_4G_PAD,IPOD_3G_PAD}{\PluginCancel}
    \opt{IPOD_4G_PAD,IPOD_3G_PAD}{\ButtonMenu}
        \opt{HAVEREMOTEKEYMAP}{& \PluginRCCancel}
        & Quit\\
\end{btnmap}

\subsubsection{Menu}
\begin{description}
    \item[Apply This Fit.] Replaces the current EQ settings with the fit
        being shown and saves them.  Unused EQ bands are disabled by setting
        their gain to 0.  If the fit has a higher peak gain than the
        original, \setting{Precut} is increased by the difference to avoid
        clipping.  To keep the original settings for later, save them with
        \setting{Save EQ Preset} before applying a fit.
    \item[Show Bands.] Lists the type, frequency, gain and Q of each band in
        the fit, and the precut that would be applied.
    \item[Revert to Original.] Restores the EQ settings that were active when
        the plugin was started.
    \item[Quit.] Exits the plugin.
\end{description}
